\section{Plasma: Woher kommt die Wärme, und wie entsteht sie?}\label{sec:plasma}

Bisher begann die Rechnung beim Plasma als Quelle von \SI{3053}{MW} (Abschnitt~\ref{sec:bilanz}). Dieser Abschnitt
fragt, wie diese Leistung entsteht und auf welchen Wegen sie das Plasma verlässt. Alle Zahlen gelten für die
HTS-Auslegung aus Stufe~0: $R_0 = \SI{10,4}{m}$, $a = \SI{2,31}{m}$, $B_0 = \SI{6,5}{T}$, Plasmavolumen
\SI{1096}{m^3}, zentrale Temperatur $T_0 = \SI{14}{keV}$. Das Skript \texttt{stellarator/plasma\_waerme.py} rechnet
jeden Schritt nach.

\subsection{Die Reaktion: Masse wird zu Energie}
Im Plasma verschmelzen Deuterium (ein Proton, ein Neutron) und Tritium (ein Proton, zwei Neutronen):
\[ \mathrm{D} + \mathrm{T} \;\rightarrow\; {}^4\mathrm{He}\;(\SI{3,5}{MeV}) + \mathrm{n}\;(\SI{14,1}{MeV}). \]
Die Produkte sind leichter als die Ausgangskerne. Mit den Atommassen \cite{ame2020} ist
\[ \Delta m = m_\mathrm{D} + m_\mathrm{T} - m_\mathrm{He} - m_\mathrm{n} = \SI{0,01888}{u}
   \quad\Rightarrow\quad Q = \Delta m\,c^2 = 0{,}01888 \cdot \SI{931,494}{MeV} = \SI{17,59}{MeV}. \]
Das sind 0,38\,\% der eingesetzten Masse. Die Energie verteilt sich nach der Impulserhaltung: Beide Produkte fliegen
mit gleich großem Impuls auseinander, also erhält das leichtere Teilchen den größeren Teil,
\[ E_\alpha = Q\,\frac{m_\mathrm{n}}{m_\mathrm{n} + m_\alpha} \approx \SI{3,54}{MeV},\qquad E_\mathrm{n} \approx \SI{14,05}{MeV}. \]
Üblich sind gerundet \SI{3,5}{} und \SI{14,1}{MeV} \cite{wiki_tritium}; die Bilanz rechnet mit \SI{3,52}{} und
\SI{14,07}{MeV}. Der Unterschied von 0,6\,\% ist für die Bilanz bedeutungslos. Wichtig ist das Verhältnis: \textbf{Vier
Fünftel der Energie trägt das Neutron, ein Fünftel der Heliumkern (Alphateilchen).} Das bestimmt alles Weitere.

\subsection{Warum es so heiß sein muss}
Beide Kerne sind positiv geladen und stoßen sich ab. Damit die Kernkraft greift, müssen sie sich bis auf die Summe ihrer
Radien nähern ($r \approx 1{,}2\,\mathrm{fm}\cdot A^{1/3}$, zusammen \SI{3,2}{fm}). Dort ist die Abstoßung
\[ E_C = \frac{e^2}{4\pi\varepsilon_0\,r} \approx \SI{444}{keV}. \]
Das Plasma hat im Zentrum \SI{14}{keV}, das sind 162~Millionen Kelvin, aber 32-mal weniger als der Wall. Dass
trotzdem Fusion stattfindet, liegt an zwei Effekten zugleich:
\begin{itemize}
\item \textbf{Maxwell-Verteilung:} Einige Teilchen sind viel schneller als der Durchschnitt; ihr Anteil fällt wie
  $e^{-E/kT}$.
\item \textbf{Tunneleffekt:} Quantenmechanisch kann ein Kern den Wall durchdringen, ohne über ihn zu steigen. Die
  Wahrscheinlichkeit steigt mit der Energie wie $e^{-\sqrt{E_G/E}}$, mit der Gamow-Energie $E_G = \SI{1182}{keV}$ für
  D-T \cite{boschhale1992}.
\end{itemize}
Das Produkt beider Kurven ist das \emph{Gamow-Fenster}. Bei \SI{14}{keV} liegt es bei
$E_0 = (\sqrt{E_G}\,kT/2)^{2/3} \approx \SI{39}{keV}$, etwa beim 2,8-Fachen von $kT$
(Bild~\ref{fig:plasma}, oben rechts). Es fusionieren also fast nur die schnellen Teilchen aus dem Ausläufer der
Verteilung.

Die Reaktionsrate fasst die Reaktivität $\langle\sigma v\rangle$ zusammen (Formel von Bosch und Hale
\cite{boschhale1992}). Sie ist bei etwa \SI{67}{keV} am größten. Für ein Kraftwerk zählt aber eine andere Frage:
Kann die Heizung durch die Alphateilchen die Verluste allein decken (Zündung)? Mit der Energieeinschlusszeit $\tau_E$
lautet die Bedingung (Lawson, ohne Strahlung)
\[ n\,T\,\tau_E \;\ge\; \frac{12\,T^2}{E_\alpha\,\langle\sigma v\rangle}, \]
und die rechte Seite ist bei \SI{13,6}{keV} am kleinsten: $\SI{2,8e21}{keV\,s\,m^{-3}}$. \textbf{Darum liegt die
Auslegungstemperatur bei \SI{14}{keV} und nicht beim Maximum der Reaktivität.} Die Auslegung erreicht im Zentrum
$n_0 T_0 \tau_E = \SI{8,0e21}{keV\,s\,m^{-3}}$. Der Wert liegt höher als die Grenze, weil Dichte und Temperatur nach
außen abfallen und die Bremsstrahlung hinzukommt.

\subsection{Wo im Plasma die Wärme entsteht}
Die Leistungsdichte ist
\[ p_\mathrm{fus}(\rho) = n_\mathrm{D}\,n_\mathrm{T}\,\langle\sigma v\rangle(T)\;Q,\qquad
   n_\mathrm{D} = n_\mathrm{T} = 0{,}425\,n_e, \]
mit den Profilen der Stufe~0: $n_e = n_0(1-\rho^2)^{0{,}5}$, $T = T_0(1-\rho^2)$, $n_0 = \SI{3,1e20}{m^{-3}}$
(Rest der Ionen: Heliumasche und Verunreinigungen). Über das Volumen integriert ergibt das \SI{3051}{MW}; die Stufe~0
nannte \SI{3053}{MW}.
\begin{itemize}
\item Im Zentrum entstehen \SI{12}{MW/m^3}, im Mittel \SI{2,8}{MW/m^3}.
\item Die Hälfte der Leistung entsteht innerhalb $\rho = 0{,}38$, also in nur 14\,\% des Volumens, 90\,\% innerhalb
  $\rho = 0{,}62$. Die Wärme entsteht im heißen Kern, der Rand trägt fast nichts bei.
\item Es laufen $1{,}08\cdot10^{21}$ Reaktionen je Sekunde. Das verbraucht am Tag \SI{0,47}{kg} Tritium und
  \SI{0,31}{kg} Deuterium. Davon werden \SI{2,9}{g} Masse in Energie umgewandelt.
\end{itemize}

\subsection{Zwei Wege aus dem Plasma}
\paragraph{Das Neutron fliegt sofort hinaus.} Es ist ungeladen, also wirkt das Magnetfeld nicht auf es. Mit 17\,\% der
Lichtgeschwindigkeit durchquert es den kleinen Radius in etwa \SI{45}{ns}, ohne das Plasma zu heizen. Seine Energie
gibt es erst in der Wand ab, vor allem im Blanket: durch Stöße mit Kernen, die es abbremsen, und durch Einfang. Die
wichtigste Einfangreaktion erzeugt zugleich neues Tritium:
\[ {}^6\mathrm{Li} + \mathrm{n} \;\rightarrow\; {}^4\mathrm{He} + \mathrm{T} + \SI{4,8}{MeV} \quad\text{\cite{wiki_tritium}}. \]
Solche Reaktionen setzen zusätzliche Energie frei. Der Faktor $M = 1{,}15$ (Annahme wie in Stufe~0) macht aus
\SI{2442}{MW} Neutronenleistung \SI{2808}{MW} Wärme im Blanket. \textbf{Gut vier Fünftel der Kraftwerkswärme
entstehen also gar nicht im Plasma, sondern im Volumen des Blankets.}

\paragraph{Das Alphateilchen bleibt und heizt.} Es ist doppelt positiv geladen und kreist um die Feldlinien, mit einem
Larmorradius von
\[ r_L = \frac{m_\alpha v_\alpha}{2e\,B_0} = \SI{4,2}{cm} \ll a = \SI{2,31}{m}. \]
So bleibt es eingeschlossen und gibt seine \SI{3,5}{MeV} über Coulomb-Stöße an das Plasma ab. Die Abbremszeit an den
Elektronen ist $\tau_s \approx \SI{0,2}{s}$ \cite{stix1972}. Oberhalb der kritischen Energie
$E_c \approx 14{,}8\,A_\alpha T_e\,(\sum n_i Z_i^2/n_e A_i)^{2/3} \approx \SI{415}{keV}$ bremsen vor allem die
Elektronen, darunter die Ionen. Über den ganzen Weg gehen daher etwa 80\,\% an die Elektronen und 20\,\% an die
Ionen. Diese \emph{Selbstheizung} von \SI{611}{MW} hält das Plasma heiß. Die Auslegung ist gezündet (keine
Fremdheizung im Betrieb), erreicht das aber nur mit dem 1,29-Fachen des Einschlusses nach der Skalierung ISS04
\cite{yamada2005}.

\paragraph{Gleichgewicht: was hereinkommt, muss hinaus.} Im stationären Betrieb gilt
\[ P_\alpha = \frac{W}{\tau_E} + P_\mathrm{Brems}:\qquad
   \SI{611}{MW} = \frac{\SI{877}{MJ}}{\SI{1,82}{s}} + \SI{129}{MW} = \SI{482}{MW} + \SI{129}{MW}. \]
\begin{itemize}
\item \textbf{Bremsstrahlung} (\SI{129}{MW}): Elektronen strahlen bei Stößen mit Ionen Röntgenlicht ab. Das trifft
  gleichmäßig die ganze Erstwand.
\item \textbf{Transport} (\SI{482}{MW}): Wärme und Teilchen wandern langsam quer zum Feld nach außen, erreichen die
  Abschälschicht (scrape-off layer) und strömen dort entlang der Feldlinien auf den Divertor.
\end{itemize}

\bild{plasma_waerme.png}{1.0}{Woher die Wärme kommt. Oben links: Reaktivität und nötiges $nT\tau_E$ für Zündung, mit
dem Auslegungspunkt. Oben rechts: Gamow-Fenster bei \SI{14}{keV}. Unten links: Profile und Leistungsdichte über dem
kleinen Radius. Unten rechts: Wege der Energie in MW.\label{fig:plasma}}
\lesehilfe{Oben links in Blau, wie häufig Deuterium und Tritium bei einer Temperatur fusionieren
(\emph{Reaktivität}). In Orange, wie gut man die Wärme einschließen muss, damit das Plasma sich selbst heizt
(\emph{Zündung}). Die Temperatur steht in keV; \SI{1}{keV} entspricht etwa 11,6 Millionen Grad. Die orange Kurve
hat ihr Minimum bei etwa \SI{14}{keV}, dort liegt der Auslegungspunkt. Oben rechts das \emph{Gamow-Fenster}: Blau
ist der Anteil der Teilchen mit einer bestimmten Energie, der steil abfällt. Orange ist die Chance, die
elektrische Abstoßung zu überwinden, die steil ansteigt. Ihr Produkt (grün) zeigt, dass die meisten Fusionen bei
etwa \SI{39}{keV} stattfinden, also bei den wenigen schnellen Teilchen, nicht beim Durchschnitt. Unten links über
dem Radius (0 Mitte, 1 Rand) Temperatur, Dichte und in Rot die erzeugte Leistung je Volumen: Die Hälfte entsteht im
innersten Siebtel des Volumens. Unten rechts, wohin die Energie geht, in MW.}


\subsection{Was daraus für die Anlage folgt}
\begin{enumerate}
\item \textbf{Die Aufteilung 80/20 ist Kernphysik, nicht Auslegung.} Das Blanket muss 82\,\% der Wärme aufnehmen, und
  zwar im Volumen, nicht an der Oberfläche. Das passt zu dem dicken, wassergekühlten Blanket nach WCLL
  (Abschnitt~\ref{sec:bilanz}).
\item \textbf{Die Wärmebilanz hat den Divertor wohl zu günstig angesetzt.} Sie teilt die Alphaleistung je zur Hälfte
  auf Erstwand und Divertor ($f_\mathrm{div} = 0{,}5$). Ohne weitere Maßnahmen gehen aber 79\,\% als Transport zum
  Divertor, also \SI{482}{MW} statt \SI{305}{MW}. Bei \SI{10}{MW/m^2} wären dann mindestens \SI{48}{m^2}
  Divertorfläche nötig statt \SI{31}{m^2}. Die Aufteilung 50/50 setzt voraus, dass etwa \SI{180}{MW} zusätzlich am
  Rand abgestrahlt werden, etwa durch gezielt eingeblasenes Fremdgas. Das ist eine Aufgabe für die Auslegung des
  Divertors.
\item \textbf{Anfahren braucht Fremdheizung.} Bis die Alphaheizung trägt, muss das Plasma von außen auf Temperatur
  gebracht werden. W7-X nutzt dafür Mikrowellen (Elektron-Zyklotron-Resonanzheizung, bis \SI{15}{MW}),
  Neutralteilcheninjektion (bis \SI{8}{MW}) und Ionen-Zyklotron-Heizung \cite{wiki_w7x}.
\item \textbf{Grenzen dieser Rechnung:} 0-D-Bilanz mit festen Profilen, nur Bremsstrahlung (keine Linien- und
  Zyklotronstrahlung), Lawson-Grenze ohne Profile, Abbremsung mit $\ln\Lambda = 17$ und Zentralwerten. Eine
  Transportrechnung, etwa mit einem 1-D-Code, wäre der nächste Schritt.
\end{enumerate}
